\begin{lemma}[Step 11: the limit is a fit of class $256$ at every centre and scale]
Let $d\in\N$ and write $\mathcal H=\ell^2(\R^d\times\R)$. Let
$\lambda:\N\to\R$, let $v_i:\mathcal H\to\R$ for each $i\in\N$, let
$\mathcal B:\R\to\mathcal P(\mathcal H)$, let $C_3:\R\times\R\to\R$, let
$i_5:\R\times\R\to\N$, let $\psi:\N\to\N$ and let $w:\mathcal H\to\R$. Assume:
\begin{enumerate}
\item\label{it:lifes-lam} $\lambda_i>0$ for every $i$, and $\lambda_i\to0$ as
$i\to\infty$;
\item\label{it:lifes-box} the boxes are the coordinate boxes,
$\mathcal B_K=\{y\in\mathcal H:|y_x|\le K\text{ and }|y_t|\le K\}$, where
$y=(y_x,y_t)$;
\item\label{it:lifes-C3} $C_3(K,S)\ge0$ for all $K,S$;
\item\label{it:lifes-psi} $\psi$ is strictly monotone;
\item\label{it:lifes-conv} $v_{\psi(j)}\to w$ uniformly on every compact subset
of $\mathcal H$ as $j\to\infty$;
\item\label{it:lifes-fit} (the uniform bound on each box) for every $K\ge1$,
every $S>0$, every $i\ge i_5(K,S)$ and every $y\in\mathcal B_K$ there are
$p'\in\R^d$ and $q'\in\R$ with
\[
  |p'|+|q'|\le C_3(K,S),\qquad \tfrac12\le1+\lambda_iq' ,
\]
such that for all $h\in\R^d$ and $\ell\in\R$ with $|\ell|\le S/2$ and
$\bigl|h-\bigl(\ell/(1+\lambda_iq')\bigr)p'\bigr|\le S$,
\[
  \bigl|v_i(y+(h,\ell))-v_i(y)-p'\cdot h-q'\ell\bigr|\le 64\,S .
\]
\end{enumerate}
Then for every centre $\hat z\in\mathcal H$ and every radius $R'>0$ there are
$p\in\R^d$ and $a\in\R$ such that $(p,a)$ is a fit for $w$ at $\hat z$ of
radius $R'$ with constant $256$, in the sense of \noderef{IsFit}.
\end{lemma}
\begin{proof}
Fix a centre $\hat z=(\hat x,\hat t)\in\mathcal H$ and a radius $R'>0$. Put
\[
  K=\max\bigl\{1,|\hat x|,|\hat t|\bigr\},\qquad S=4R' ,
\]
so $K\ge1$ and $S>0$, and $\hat z\in\mathcal B_K$ by \ref{it:lifes-box},
because $|\hat x|\le K$ and $|\hat t|\le K$ by the definition of $K$.

\emph{The per-index coefficients.} For every $j$ with $\psi(j)\ge i_5(K,S)$,
hypothesis \ref{it:lifes-fit} applied at the index $i=\psi(j)$ and the point
$y=\hat z$ supplies $p'_j\in\R^d$ and $q'_j\in\R$ with
\begin{equation}\label{eq:lifes-bd}
  |p'_j|+|q'_j|\le C_3(K,S),\qquad \tfrac12\le A_j:=1+\lambda_{\psi(j)}q'_j ,
\end{equation}
and with the increment bound of \ref{it:lifes-fit}. Since $\psi$ is strictly
monotone by \ref{it:lifes-psi} we have $\psi(j)\ge j$, so all $j\ge i_5(K,S)$
qualify; discard the finitely many smaller indices.

\emph{Extraction.} By \eqref{eq:lifes-bd} the sequences $(p'_j)$ and $(q'_j)$
are bounded, the first in $\R^d$ and the second in $\R$, so by
Bolzano--Weierstrass --- available because $\R^d$ is finite-dimensional, hence
its closed balls are compact --- there are $p\in\R^d$, $q\in\R$ and a strictly
monotone $l\mapsto j_l$ with $p'_{j_l}\to p$ and $q'_{j_l}\to q$. Because
$\lambda_i\to0$ by \ref{it:lifes-lam} and $|q'_j|\le C_3(K,S)$ uniformly,
$\lambda_{\psi(j)}q'_j\to0$ and therefore
\begin{equation}\label{eq:lifes-A}
  A_{j_l}\to1,\qquad\text{hence}\qquad \frac{p'_{j_l}}{A_{j_l}}\to p ,
\end{equation}
the quotient being legitimate throughout since $A_j\ge\tfrac12>0$.

\emph{The admissible increments.} Let $h\in\R^d$ and $\ell\in\R$ satisfy
$|\ell|\le R'$ and $|h-\ell p|\le R'$. Then $|\ell|\le R'=S/4\le S/2$, and by
the triangle inequality
\[
  \Bigl|h-\frac{p'_{j_l}}{A_{j_l}}\,\ell\Bigr|
   \le|h-\ell p|+|\ell|\,\Bigl|p-\frac{p'_{j_l}}{A_{j_l}}\Bigr|
   \le R'+R'\Bigl|p-\frac{p'_{j_l}}{A_{j_l}}\Bigr| .
\]
By \eqref{eq:lifes-A} the last factor is at most $1$ for all large $l$, and the
threshold depends only on the sequence, not on $(h,\ell)$; for such $l$ the
left side is at most $2R'\le S$. Hypothesis \ref{it:lifes-fit} therefore applies
at the index $\psi(j_l)$, the centre $\hat z\in\mathcal B_K$ and this pair
$(h,\ell)$, giving
\begin{equation}\label{eq:lifes-pre}
  \bigl|v_{\psi(j_l)}(\hat z+(h,\ell))-v_{\psi(j_l)}(\hat z)
    -p'_{j_l}\cdot h-q'_{j_l}\ell\bigr|\le 64\,S=256\,R' .
\end{equation}

\emph{Passing to the limit.} The set
\[
  \mathcal K=\bigl\{\hat z+(h,\ell):|\ell|\le R',\ |h-\ell p|\le R'\bigr\}
\]
is closed and bounded in $\mathcal H$, hence compact: boundedness holds because
$|\ell|\le R'$ and $|h|\le|h-\ell p|+|\ell||p|\le R'(1+|p|)$. By
\ref{it:lifes-conv}, $v_{\psi(j_l)}\to w$ uniformly on $\mathcal K$, and in
particular pointwise at the two points $\hat z+(h,\ell)$ and $\hat z$, both of
which lie in $\mathcal K$ (the second by taking $h=0$, $\ell=0$). Since also
$p'_{j_l}\to p$ and $q'_{j_l}\to q$, the quantity inside the absolute value in
\eqref{eq:lifes-pre} converges to
$w(\hat z+(h,\ell))-w(\hat z)-p\cdot h-q\ell$, and a non-strict inequality is
preserved in the limit:
\begin{equation}\label{eq:lifes-post}
  \bigl|w(\hat z+(h,\ell))-w(\hat z)-p\cdot h-q\ell\bigr|\le256\,R'
\end{equation}
whenever $|\ell|\le R'$ and $|h-\ell p|\le R'$.

\emph{Conclusion.} Put $a=q+|p|^2/2$, so that $q=a-|p|^2/2$. Then
\eqref{eq:lifes-post} is exactly the assertion that $(p,a)$ is a fit for $w$ at
$\hat z$ of radius $R'$ with constant $256$ in the sense of \noderef{IsFit}.
As $\hat z$ and $R'$ were arbitrary, this proves the lemma. The auxiliary
subsequence $l\mapsto j_l$ depends on $\hat z$ and $R'$, which is harmless:
\noderef{IsFit} is asserted separately for each centre and each radius, and
$w$ itself does not depend on the extraction.
\end{proof}
